\begin{lemma}
\label{lemma-integral-over-field}
Let $k$ be a field. Let $S$ be a $k$-algebra over $k$.
\begin{enumerate}
\item If $S$ is a domain and finite dimensional over $k$,
then $S$ is a field.
\item If $S$ is integral over $k$ and a domain,
then $S$ is a field.
\item If $S$ is integral over $k$ then every prime of
$S$ is a maximal ideal (see
Lemma \ref{lemma-ring-with-only-minimal-primes}
for more consequences).
\end{enumerate}
\end{lemma}

\begin{proof}
For (3), a quotient $S/\mathfrak q$ by any prime is an integral
$k$-algebra and a domain. Its coefficient map from $k$ is injective,
since the quotient is nonzero and $k$ is a field. Assertion (2)
makes this quotient a field, hence $\mathfrak q$ is maximal.
Let $S$ be integral over $k$ and assume $S$ is a domain.
Take a nonzero $s\in S$. By Lemma
\ref{lemma-characterize-integral} we may find a
finite dimensional $k$-subalgebra $k \subset S' \subset S$
containing $s$. This subalgebra is a domain since it is a subring
of $S$. Once (1) is proved, $S'$ is a field and the inverse of
the original nonzero $s$ belongs to $S'\subset S$, proving (2).
For (1), assume that the original $S$ is finite dimensional over
$k$ and is a domain. Pick a nonzero $s\in S$. The $k$-linear
map $S\to S$, $a\mapsto sa$, is injective because $S$ is a
domain and $s\ne0$. Equality of the finite source and target
dimensions makes it surjective. A preimage $s_1$ of $1$ satisfies
$ss_1=1$. Thus every nonzero element is invertible and $S$ is a field.
\end{proof}